% TOP_PROBLEM: {"id": 11, "title": "Twin Prime Conjecture", "category_id": 1, "category": {"id": 1, "name": "number_theory", "display_name": "Number Theory", "description": "Properties of integers, prime numbers, Diophantine equations.", "slug": "number-theory", "order_index": 1, "created_at": "2026-07-31T15:26:25.670Z"}, "status": "open"}
% FIRST_POSTED: 2026-09-27
% !TeX program = lualatex
% ATTEMPT_STATUS: unresolved
% Research continuation by GPT_6_Astra_Ultra, 27 September 2026.
% Catalogue statements and existing sources preserved verbatim.
\documentclass[11pt]{article}
\usepackage[margin=25mm]{geometry}
\usepackage{fontspec}
\setmainfont{Latin Modern Roman}
\usepackage{amsmath,amssymb,mathtools}
\usepackage{unicode-math}
\setmathfont{Latin Modern Math}
\usepackage{xurl,hyperref}
\hypersetup{colorlinks=true,urlcolor=blue}
\setcounter{secnumdepth}{0}
\setlength{\parindent}{0pt}
\setlength{\parskip}{5pt}
\setlength{\emergencystretch}{3em}
\begin{document}
\section*{\texorpdfstring{Twin Prime Conjecture}{Twin Prime Conjecture}}\label{11}
\subsection{Definitions and mathematical statement}
Let $\mathcal P=\{p\in{\mathbb{Z}}_{\ge2}:p\text{ has exactly two positive divisors}\}$. The conjecture is
\[
 \forall X>0\ \exists p>X:\qquad p\in\mathcal P,\quad p+2\in\mathcal P.
\]
Equivalently, $\#\{p\le x:p,p+2\in\mathcal P\}\to\infty$. No particular asymptotic formula for this counting function is required.
\par\smallskip\textit{Formulation and concept references:} \href{https://mathworld.wolfram.com/TwinPrimeConjecture.html}{[S1]}.

\subsection{Short English statement}
Are prime pairs separated by exactly two found arbitrarily far along the integers?
\subsection{Research attempt}

\paragraph{Target and current context.}
Only infinitude of the gap two is required; a Hardy--Littlewood
asymptotic would be stronger. The established Polymath result [E1]
gives $\liminf(p_{n+1}-p_n)\le246$. Stadlmann's August 2026 manuscript
[E2] states the improvement to $240$; its statement and construction were
inspected, but its full proof was not independently audited here. Neither
bound identifies the recurring gap as two. From infinitely many gaps in
a finite even set one can select some recurring gap by pigeonhole, but
cannot prescribe which member of that set recurs.

A second established approximation is Chen's theorem: infinitely many
primes $p$ have $p+2$ with at most two prime factors, counted with
multiplicity. Green--Tao [E3], Section 6, records a quantitative version
in dyadic intervals and develops further consequences for these primes.
The present attempt starts from that approximation and isolates exactly
what would distinguish its prime outputs from semiprime outputs. The
catalogue's claimed proofs [S3]--[S5] are retained as sources, not assumed
as established resolutions. No result below relies on them.

\paragraph{An exact parity reduction on Chen candidates.}
Let $\Omega(n)=\sum_{p^a\parallel n}a$ and
$\lambda(n)=(-1)^{\Omega(n)}$. Define
\[
 \mathcal C(x)=\{p\le x:p\text{ prime},\ \Omega(p+2)\le2\},\qquad
 C(x)=|\mathcal C(x)|,\qquad
 L_C(x)=\sum_{p\in\mathcal C(x)}\lambda(p+2).
\]
Since $p+2\ge4$, there is no zero-factor case. On this particular set,
$\lambda(p+2)=-1$ exactly when $p+2$ is prime, and it equals $+1$
exactly when $p+2$ is a product of two primes, including a square.
Consequently the twin-prime counting function satisfies the exact identity
\begin{equation}\label{ra:20:parity}
 \pi_2(x)=\frac{C(x)-L_C(x)}2.
\end{equation}
No probabilistic independence or asymptotic estimate is involved.

\textbf{Conditional consequence.} If for some fixed $\delta>0$ and
an unbounded sequence $x_j$ one could prove
\begin{equation}\label{ra:20:gap}
 L_C(x_j)\le(1-\delta)C(x_j),
\end{equation}
then the twin-prime conjecture would follow. Indeed Chen's theorem gives
$C(x)\to\infty$, and \eqref{ra:20:parity} gives
$\pi_2(x_j)\ge\delta C(x_j)/2\to\infty$.
For the stated quantitative version of Chen's theorem, the same condition
would give $\pi_2(x_j)\gg_\delta x_j/(\log x_j)^2$.
Only the infinitude conclusion is needed; the stronger density statement
is not silently built into the original problem.

The trivial estimate is $L_C(x)\le C(x)$, which loses everything at
the endpoint. A proof that the candidate set is large cannot by itself
improve that estimate: it is logically compatible with every sufficiently
large candidate being a semiprime. Full cancellation $L_C=o(C)$ would
be more than enough, but the positive saving in \eqref{ra:20:gap} is
already sufficient. Nothing in this notebook proves that saving.

There is a useful distinction between the exact reformulation and the
chosen sufficient estimate. The original conjecture is equivalent to
$C(x)-L_C(x)\to\infty$, by the identity and monotonicity of the twin
count. It does not require that this deficit be a fixed proportion of
$C(x)$. A varying saving $\delta(x)>0$ would also suffice if
$\delta(x)C(x)\to\infty$ and $L_C(x)\le(1-\delta(x))C(x)$ along
unbounded $x$. Thus failure to prove a fixed $\delta$ would not refute
the original problem. The restricted sign sum is being used as a
concrete attack target with sufficient quantitative strength, while the
weaker infinitude quantifier remains the actual objective.

\paragraph{Why ordinary parity information does not supply it.}
The word parity here concerns the number of prime factors, not whether
$p$ or $p+2$ is odd. All sufficiently large candidates already consist
of odd numbers. It is also essential that $\Omega$, rather than the
number of distinct prime factors, is used: $9$ has two prime factors
with multiplicity and must receive sign $+1$.

A global mean-value statement for $\lambda(n)$ over all integers would
not automatically control its mean on the much thinner set
$\{p+2:p\in\mathcal C(x)\}$. As a purely finite illustration, any
balanced list of $+1$ and $-1$ values has mean zero, while its subset of
positive values has mean one. This is not a substitute model for the
actual Liouville function; it shows why a mean estimate must survive the
specific primality and almost-primality restrictions in the displayed
sum. Treating those restrictions as independent of the sign is precisely
an unproved correlation assumption.

Polymath's discussion of the sieve parity obstruction [E1, Section 8]
explains the limitation for the class of sieve methods considered there.
It is not an impossibility theorem for all number-theoretic methods.
Here \eqref{ra:20:parity} identifies the missing distinction explicitly,
without claiming that merely naming the parity barrier solves it.

\paragraph{A second concrete target: weighted two-point correlation.}
Define $\Lambda(p^a)=\log p$ for prime powers with $a\ge1$, and
$\Lambda(n)=0$ otherwise. Put
\[
 T_2(x)=\sum_{2\le n\le x}\Lambda(n)\Lambda(n+2),\qquad
 P_2(x)=\sum_{\substack{p\le x\\p,p+2\text{ prime}}}
                                      \log p\log(p+2).
\]
The weighted sum is convenient algebraically, but prime powers must be
removed before a positive value can certify a twin pair.

\textbf{Lemma 20.1.} For $x\ge2$,
\[
 0\le T_2(x)-P_2(x)
 \le2(\log(x+2))^2
       \sum_{a=2}^{\lfloor\log_2(x+2)\rfloor}(x+2)^{1/a}
       =O(\sqrt x(\log(x+2))^3).
\]
\textit{Proof.}
Every nonzero term absent from $P_2$ has at least one coordinate equal
to a proper prime power. The number of such powers at most $x+2$ is
at most the displayed sum, since counting every integer base only
increases it. Either coordinate supplies at most that many indices
$n$, hence the factor two. Each weight is at most $(\log(x+2))^2$.
Finally the square terms contribute at most $\sqrt{x+2}$ and the
remaining $O(\log(x+2))$ exponents contribute at most $(x+2)^{1/3}$
each. This gives the stated bound, with no assumption about primes.
\hfill$\square$

Therefore
\begin{equation}\label{ra:20:weighted-lower}
 \pi_2(x)\ge
 \frac{T_2(x)-O(\sqrt x\log^3(x+2))}{\log^2(x+2)}.
\end{equation}
A lower bound $T_2(x_j)\ge c x_j$ for some fixed $c>0$ along an
unbounded sequence would prove infinitude. In fact it suffices that
$T_2(x_j)/(\sqrt{x_j}\log^3(x_j+2))\to\infty$.
These are sufficient conditions, not asserted equivalences: the original
conjecture permits a much sparser infinite family. For small $x$,
$T_2(x)>0$ alone proves nothing; the pair $(2,4)$ contributes before
any twin pair is counted, and $(25,27)$ consists entirely of proper
prime powers.

\paragraph{Attempt to evaluate the correlation by divisor identities.}
The exact identity
\[
 \Lambda(n)=-\sum_{d\mid n}\mu(d)\log d
\]
follows by expanding over the distinct prime divisors of $n$. If there
is exactly one distinct prime, only $d=1,p$ contribute; if there are
at least two, the coefficient of each $\log p$ is
$-\prod_{q\mid n,\ q\ne p}(1-1)=0$.
Thus, writing $A_{d,e}(x)$ for the number of $2\le n\le x$ with
$d\mid n$ and $e\mid n+2$, we get
\[
 T_2(x)=\sum_{d\le x}\sum_{e\le x+2}
                  \mu(d)\mu(e)\log d\log e\,A_{d,e}(x).
\]
The two congruences are compatible exactly when $\gcd(d,e)\mid2$.
If compatible, the solutions form one residue class modulo $[d,e]$,
so $A_{d,e}(x)=x/[d,e]+O(1)$; otherwise the count is zero.

This exposes explicit main terms but does not justify discarding the
errors. Summing their absolute values over the full divisor rectangle
gives only $O(x^2\log^2(x+2))$, far larger than the desired linear
scale. Truncating to short divisors reduces the summation error but no
longer equals $\Lambda$. For example, if one retains only divisors
$d\le R$ in the displayed formula, a prime $p>R$ receives zero instead
of $\log p$. Any repaired truncation requires an analysis of the omitted
terms and their interaction between $n$ and $n+2$.
This is the first explicit failure of the correlation calculation: local
solvability of congruences is easy, but summing the signed errors and
long-divisor terms at the required scale is not accomplished.

\paragraph{Bounded gaps, fixed gaps, and false extrapolation.}
An infinite increasing integer sequence with every successive gap four
has bounded gaps and no gap two. This auxiliary example does not model
all properties of primes, but refutes the bare implication from a
finite gap bound to twins. Stronger results finding two primes somewhere
inside an admissible tuple have the same indexing issue: the successful
pair may vary. Pigeonhole yields a fixed pair recurring infinitely often,
but not necessarily the pair with offsets zero and two.

Nor does increasing the number of computed twins settle the quantifier
$\forall X\ \exists p>X$. Every finite prefix can be extended as an
abstract prime-like list with no further gap two. Such a list is not an
alternative set of actual primes; it only isolates why a finite table
requires a structural argument controlling the tail. Our two candidate
estimates, \eqref{ra:20:gap} and \eqref{ra:20:weighted-lower}, specify
the missing tail information rather than inferring it from the table.

\paragraph{Finite verification and outcome.}
Exact factorization by a small-prime sieve gives
\[
\begin{array}{r|r|r|r|r}
x&C(x)&L_C(x)&\pi_2(x)&\text{proper-prime-power terms}\\\hline
100&20&4&8&9\\1000&115&45&35&15\\
10000&633&223&205&24
\end{array}
\]
The final column counts indices contributing to $T_2-P_2$, not their
logarithmic weights. The script verifies \eqref{ra:20:parity} exactly
and bounds the number of these indices using exact integer roots.
It also evaluates the logarithmic sums as ordinary floating-point
diagnostics; at $x=10000$ they are approximately $13302.9243$ and
$12950.5862$. These decimals are not interval certificates or evidence
of an asymptotic lower bound. Code and output are in
\texttt{\_Local/top\_problems/continuation\_15\_25\_2026\_09\_27/analytic/}.

\textbf{Outcome: UNRESOLVED.} The checked partial work is an exact
prime-versus-semiprime identity, a proved prime-power removal bound, and
an explicit analysis of the failed divisor-correlation approximation.
The missing theorem is a positive saving in the restricted Liouville
sum, or a sufficiently large lower bound for the actual two-point
Mangoldt correlation. No such estimate, no infinite construction of
twin pairs, and no counterexample have been obtained. The computations
and elementary proofs were checked locally; the open conjecture is not
claimed solved or independently certified.

\subsection{Sources}
\begin{itemize}
\item[S1]Weisstein, E. W., Twin Prime Conjecture, From MathWorld - A Wolfram Resource, updated 16 August 2026.
\url{https://mathworld.wolfram.com/TwinPrimeConjecture.html}.
\textit{Formulation source.}
\item[S2]Elementary Number Theory — Henrik Bachmann, version 4, December 1, 2025.
\url{https://www.henrikbachmann.com/uploads/7/7/6/3/77634444/ent_2025_overviewnotes_v4.pdf}.
\textit{Further reference; not a proof endorsement.}
\item[S3]A New Look on the Twin-Prime Conjecture: The Proof of the Twin-Prime Conjecture — Woodward and Asheralieva.
\url{https://researchportal.hw.ac.uk/en/publications/a-new-look-on-the-twin-prime-conjecture-the-proof-of-the-twin-pri/}.
\textit{Further reference; not a proof endorsement.}
\item[S4]Proof of the Hardy-Littlewood K-tuple Conjecture in the Distribution of Numbers Coprime with the Primorial — Tim Samshuijzen, May 2026 revision.
\url{https://rxiv.org/pdf/2501.0157v4.pdf}.
\textit{Further reference; not a proof endorsement.}
\item[S5]Twin Prime Conjecture 3rd Way — Taha Muhammad, June 19, 2026.
\url{https://www.cambridge.org/engage/coe/article-details/6a2d20d8d1922e37d572c5f8}.
\textit{Further reference; not a proof endorsement.}

\item[E1]D. H. J. Polymath, \emph{Variants of the Selberg sieve, and
bounded intervals containing many primes}, Research in the Mathematical
Sciences 1 (2014), article 12; arXiv:1407.4897, especially Section 8.
\url{https://arxiv.org/pdf/1407.4897}.
\textit{Primary research source for the established bound 246 and the
specified sieve parity obstruction. Consulted 27 September 2026.}
\item[E2]Julia Stadlmann, \emph{Bounded gaps between primes},
arXiv:2608.31126v1 (31 August 2026).
\url{https://arxiv.org/html/2608.31126v1}.
\textit{Recent manuscript stating the bound 240. Statement inspected;
full proof not independently audited and not used as a premise here.
Consulted 27 September 2026.}
\item[E3]Ben Green and Terence Tao,
\emph{Restriction theory of the Selberg sieve, with applications},
Journal de Th\'eorie des Nombres de Bordeaux 18 (2006), 147--182;
arXiv:math/0405581, Section 6, Theorem 6.1.
\url{https://arxiv.org/pdf/math/0405581}.
\textit{Primary research source recording Chen's theorem and a quantitative
dyadic formulation, with additional restrictions that only strengthen the
lower bound for the larger candidate set used here.
Consulted 27 September 2026.}
\end{itemize}
\end{document}
